\section{Scattering matrices and analytic estimates}
\label{mat:section}

All matrices in this section are finite dimensional. A port at a physical
site $x$, of strength $\sigma>0$, is represented by the column
$\sqrt{\sigma}\,e_x$. In the matrix $Y$ of these columns, rows index
physical sites and columns index ports; several distinct ports may be
attached to the same site. For a real symmetric matrix
$H^\#$ and a real energy $E$, set
\begin{equation}
  \mathscr L=H^\#-E-iYY^\top,\qquad
  G=\mathscr L^{-1},\qquad
  S=\one+2iY^\top GY.
  \label{mat:response}
\end{equation}
The entry $S_{ja}$ is the scattering amplitude from input port $a$ to
output port $j$. The physical matrix $H^\#$ may differ from the true
Hamiltonian by real diagonal changes at sites carrying ports. We call these changes
\emph{reference potentials}. They may depend on the true potentials.

Absorption bounds the response between ports even when the full inverse
$G$ is large. Closing ports removes their absorptions and changes real
diagonal values; we express this operation as a rational family of
scattering matrices and compare its energies after averaging the closing
values.

\subsection{Invertibility, absorption, and Schur complements}

The inverses in \Cref{mat:response} exist almost surely under any
absolutely continuous joint law of the true site potentials. Indeed, if
$\mathscr L u=0$, taking the imaginary part of $u^*\mathscr L u$ gives
$Y^\top u=0$. Thus $u$ vanishes at every port site, so restoring the true
diagonal entries leaves its equation unchanged: $(H-E)u=0$. For a fixed
finite site set, $\det(H-E)$ is a polynomial in its true diagonal
variables whose coefficient of their full product is one. It is not
identically zero. A nonzero polynomial has a Lebesgue-null zero set,
as follows by induction on the number of variables and one-variable
slicing. This proves the assertion even when the port set and reference
potentials depend on those variables. The same argument applies to
deterministic restrictions to physical sites. We always discard the
resulting null exceptional sets when forming such inverses.

A matrix $D$ is \emph{passive} if $\operatorname{Im}D\le0$, with the
Hermitian imaginary part fixed in \cref{setup:conventions}. The absorption
identities below apply to any real rectangular $Y$, including real
coordinate or subspace inclusions.

\begin{lemma}[Absorption identities]\label{mat:ward}
Suppose $\mathscr L$ is invertible and
$W=-\operatorname{Im}\mathscr L\ge YY^\top$. Then
\begin{align}
  \operatorname{Im}G&=G^*WG=GWG^*,
  \label{mat:ward-identity}\\
  \norm{Y^\top GYa}^2
    &\le \operatorname{Im}(a^*Y^\top GYa),
  \qquad \norm{Y^\top GY}\le1.
  \label{mat:ward-contraction}
\end{align}
For the physical system in \Cref{mat:response}, $S$ is symmetric and
unitary. If $\mathscr L$ is complex symmetric and passive, each defined
Schur complement is complex symmetric and passive. An absorption lower
bound supported on retained coordinates persists in the Schur complement.
\end{lemma}

\begin{proof}
The inverse difference identities give
\[
  G-G^*=G^*(\mathscr L^*-\mathscr L)G
        =G(\mathscr L^*-\mathscr L)G^*,
\]
which proves \Cref{mat:ward-identity}. Consequently, with $K=Y^\top GY$,
\[
 \norm{Ka}^2
 \le a^*Y^\top G^*WGY a
 =\operatorname{Im}(a^*Ka)
 \le\norm a\,\norm{Ka}.
\]
This proves \Cref{mat:ward-contraction}. In the physical case $W=YY^\top$,
we have $\operatorname{Im}K=K^*K$ and hence
\[
 (\one+2iK)^*(\one+2iK)
   =\one-4\operatorname{Im}K+4K^*K=\one.
\]
Since the matrix is square, it is unitary. Symmetry follows from
$\mathscr L^\top=\mathscr L$ and the fact that $Y$ is real.

For the Schur assertion, write a symmetric matrix in blocks as
\[
 D=\begin{pmatrix}D_{11}&C\\ C^\top&D_{22}\end{pmatrix},\qquad
 D_{\mathrm{eff}}=D_{11}-CD_{22}^{-1}C^\top.
\]
For a vector $v$ on the retained coordinates, extend it to
$\widetilde v=(v,-D_{22}^{-1}C^\top v)$. Then
\[
 D\widetilde v=(D_{\mathrm{eff}}v,0),\qquad
 \widetilde v^*D\widetilde v=v^*D_{\mathrm{eff}}v.
\]
Taking imaginary parts proves passivity. If
$-\operatorname{Im}D\ge\diag(W_0,0)$, the same identity gives
$-\operatorname{Im}D_{\mathrm{eff}}\ge W_0$. Symmetry is immediate from
the displayed Schur formula.
\end{proof}

We will also use the elementary kernel consequence of passivity: if
$Du=0$ and $D$ is passive, then
$\operatorname{Im}D\,u=\operatorname{Re}D\,u=0$. Indeed, the nonnegative
matrix $-\operatorname{Im}D$ has zero quadratic form on $u$, so it
annihilates $u$; the assertion about the real part follows from $Du=0$.

\subsection{Closing ports and comparing energies}

Symmetry and unitarity make $\abs{S_{ja}}^2$ symmetric nonnegative
weights with unit row sums. For real scalar or real Euclidean-vector
data $f$ on the ports, their Dirichlet energy and its single-port
contribution are
\begin{equation}
 \mathcal Q_S(f)=\frac12\sum_{a,j}\abs{S_{ja}}^2\abs{f_j-f_a}^2,
 \qquad
 \mathcal M_a(f)=\sum_j\abs{S_{ja}}^2\abs{f_j-f_a}^2.
 \label{mat:energies}
\end{equation}
For scalar data, writing $F=\diag(f)$, unitarity gives
\begin{equation}
 \mathcal Q_S(f)=\tr(F^2-FSFS^*).
 \label{mat:energy-trace}
\end{equation}
For vector data, one sums this identity over real components.

\begin{lemma}[Common-phase closure]\label{mat:closure}
Partition the ports into retained ports and a set $C$ to be closed, and
write
\[
 S=\begin{pmatrix}A&T\\T^\top&B\end{pmatrix}.
\]
Close every port in $C$ and add the real shift $w\sigma_j$ to the
physical diagonal for each closed port $j$ of strength $\sigma_j$.
Where the resulting inverse exists, its scattering matrix on the
retained ports is
\begin{equation}
 S_z=A+Tz(\one-zB)^{-1}T^\top,
 \qquad z=\frac{w+i}{w-i}.
 \label{mat:closure-formula}
\end{equation}
The right-hand side is analytic for $\abs z<1$.
\end{lemma}

\begin{proof}
If $Y_C$ consists of the closed port columns, the site matrix becomes
$\mathscr L_w=\mathscr L+(w+i)Y_CY_C^\top$. Woodbury inversion and
$Y_C^\top GY_C=(B-\one)/(2i)$ give the result, using
\[
 \frac{z}{1-z}=-\frac{w+i}{2i},\qquad
 \one+(w+i)Y_C^\top GY_C=\frac{\one-zB}{1-z}.
\]
These are rational identities, so the formulas extend across any
removable singularities. The determinant of $\mathscr L_w$ is a
polynomial in $w$ which is nonzero at $w=-i$; hence only finitely many
closing values can make the site matrix singular. Finally $B$ is a
compression of a unitary matrix and is therefore contractive, which
makes $\one-zB$ invertible in the open disk.
\end{proof}

\begin{lemma}[Conditional single-port comparison]
\label{mat:energy-closure}
Consider closing a single port $a$. Set
\[
 u=S_{\widehat a,a},\qquad b=S_{aa},\qquad
 \gamma=\frac{z}{1-zb},
\]
where $\widehat a$ means that port $a$ is omitted. Let $\mathcal F$
contain the open system and the real data $f$, before the closing value
is sampled. Assume that its conditional law assigns no mass to the
finite exceptional set in \Cref{mat:closure}. If $\abs b<1$, put
$\mu=\E[\gamma\mid\mathcal F]$. Then
\begin{equation}
 \E[\mathcal Q_{S_z}(f_{\widehat a})\mid\mathcal F]
 \le \mathcal Q_S(f)+2\abs\mu\,\mathcal M_a(f).
 \label{mat:conditional-comparison}
\end{equation}
If $\abs b=1$, the same assertion holds with zero error, and we use the
convention $\mu=0$.
\end{lemma}

\begin{proof}
First take scalar data and subtract $f_a$ from every datum. This does
not change either energy, and it makes the removed datum zero. Let
$F=\diag(f_{\widehat a})$ after this subtraction and
$A=S_{\widehat a,\widehat a}$. By \Cref{mat:energy-trace}, expansion of
$S_z=A+\gamma uu^\top$ yields the exact identity
\[
 \mathcal Q_{S_z}(f_{\widehat a})-\mathcal Q_S(f)
  =-2\operatorname{Re}\bigl(\gamma\,u^\top F A^*F u\bigr)
    -\abs\gamma^2\,(u^*Fu)^2.
\]
The last square is real and nonnegative. Moreover,
\[
 \abs{u^\top F A^*Fu}\le\norm{Fu}^2=\mathcal M_a(f),
\]
because $A$ is contractive. Conditional expectation therefore proves
\Cref{mat:conditional-comparison}. For vector data the same argument is
applied componentwise. If $\abs b=1$, unitarity gives $u=0$, so closing
the port leaves the scattering matrix on the other ports unchanged.
\end{proof}

The use of the conditional mean $\mu$, rather than
$\E[\abs\gamma\mid\mathcal F]$, is important. In later applications the
open system and the data must remain measurable without exposing the
hidden true value at the port to be closed.

A standard Cauchy law for $w$ makes this conditional mean vanish. To
restore a potential within a bounded interval, we truncate that law and
estimate the resulting error.

\begin{lemma}[A cutoff-Cauchy closing law]\label{mat:cutoff}
Fix a smooth function $\chi:\R\to[0,1]$ equal to one near zero and
supported in a fixed sufficiently small interval about zero. For a
port of strength $\sigma$, give its hidden real shift $\delta$ the law
\begin{equation}
 \frac1{Z_\sigma}\chi(\delta)
       \frac{\sigma\,\dd\delta}{\pi(\delta^2+\sigma^2)},
 \qquad
 Z_\sigma=\int_\R\chi(\delta)
       \frac{\sigma\,\dd\delta}{\pi(\delta^2+\sigma^2)}.
 \label{mat:cutoff-law}
\end{equation}
Use $w=\delta/\sigma$ in \Cref{mat:closure}. Conditional on the open
system, let $g=\sigma G_{xx}$, where $x$ is the site of the port. For all
sufficiently small $\sigma$, the mean in
\Cref{mat:energy-closure} satisfies
\begin{equation}
 \abs\mu\le C,\qquad
 \abs\mu\le C\sqrt\sigma\quad\hbox{if }\abs g>\sqrt\sigma.
 \label{mat:cutoff-estimate}
\end{equation}
The constants depend on $\chi$, not on the exterior system.
\end{lemma}

\begin{proof}
The case $\abs{S_{aa}}=1$ uses the zero-error convention. Otherwise put
$b=S_{aa}=1+2ig$, so $\abs b<1$. Under the uncut Cauchy law,
\[
 z=\frac{\delta+i\sigma}{\delta-i\sigma}=e^{it},\qquad
 \delta=\sigma\cot(t/2),
\]
is Haar distributed on the circle. Its mean of
$z/(1-bz)=\sum_{k\ge0}b^kz^{k+1}$ is zero. Define the deleted angular
density
\[
 d_\sigma(t)=1-\chi\bigl(\sigma\cot(t/2)\bigr).
\]
It is bounded by one, supported on an arc of length $O(\sigma)$ about
$t=0$, and Lipschitz with constant $O(\sigma^{-1})$. To see the last
claim, derivatives can occur only when $\delta$ lies in a fixed compact
set separated from zero; there $\abs t$ is comparable to $\sigma$ and
$\abs{\dd\delta/\dd t}=O(\sigma^{-1})$. Near $t=0$ the function is
constant. Thus $Z_\sigma=1-O(\sigma)$ and
\begin{equation}
 \mu=-Z_\sigma^{-1}\int_\T
           d_\sigma(t)\frac{e^{it}}{1-be^{it}}\,\dd m_\T(t),
 \label{mat:deleted-density}
\end{equation}
where $\dd m_\T(t)=\dd t/(2\pi)$.

On the support of $d_\sigma$,
\[
 \abs{1-be^{it}}\ge\abs{1-b}-\abs b\abs{1-e^{it}}
                    \ge 2\abs g-C\sigma.
\]
If $\abs g>\sqrt\sigma$, this is at least $\abs g$ for small enough
$\sigma$. The integral in \Cref{mat:deleted-density} is then bounded
by $C\sigma/\abs g\le C\sqrt\sigma$.

For the uniform bound, $\abs b\le1/2$ is immediate. If
$1/2<\abs b<1$, use
\[
 \frac z{1-bz}=b^{-1}\left(\frac1{1-bz}-1\right).
\]
Write $b=\rho e^{i\beta}$ and use
\[
 \frac1{1-\rho e^{it}}
   =\frac12\bigl(1+P_\rho(t)+iQ_\rho(t)\bigr),
\]
where
\[
 P_\rho(t)=\frac{1-\rho^2}{1-2\rho\cos t+\rho^2},\qquad
 Q_\rho(t)=\frac{2\rho\sin t}{1-2\rho\cos t+\rho^2}.
\]
The Poisson contribution is bounded since $P_\rho\ge0$ has integral
one. For the conjugate contribution, if the singular center $-\beta$
is farther than a sufficiently large multiple of $\sigma$ from the
support arc, use $\abs{Q_\rho(t)}\le C/\abs t$ in angular distance to
bound the integral absolutely. Otherwise, after recentering at
$-\beta$, the density is supported in $\abs t\le C'\sigma$ and has
Lipschitz constant $C/\sigma$. Oddness gives
\[
 \abs{\int Q_\rho(t)d_\sigma(t-\beta)\,\dd m_\T(t)}
 \le C\int_0^{C'\sigma}\frac1t\frac t\sigma\,\dd t\le C.
\]
The constant part and normalization in
\Cref{mat:deleted-density} are harmless. This proves both bounds.
\end{proof}

\subsection{Dimension-independent analytic bounds}

The common-phase closure formula contains the resolvent
$(\one-zB)^{-1}$. We need weak-$L^1$ bounds for its sandwiched and trace
expressions, with no dimensional factor beyond the displayed norms.
For a nonnegative random variable $X$ on the circle, write
\[
 \norm X_{L^{1,\infty}}
    =\sup_{t>0}t\,m_\T\{X>t\}.
\]
The Hilbert--Schmidt estimate below has the same dimension-independent
scale as the dissipative-resolvent bound in
\cite[Lemma~3.1 and Appendix~C]{AizenmanElgartNabokoSchenkerStolz2006}.
We prove the disk estimates directly using a Hilbert-space-valued
Calder\'on--Zygmund decomposition
\cite[Theorem~1]{BenedekCalderonPanzone1962}.

\begin{lemma}[Weak-$L^1$ resolvent estimate]\label{mat:weak}
Let $B$ be a contraction and let $U,V$ be compatible deterministic
matrices. At the almost-everywhere boundary values of the rational
function in question,
\begin{equation}
 \left\|\norm{U(\one-zB)^{-1}V}_{\mathrm{HS}}\right\|_{L^{1,\infty}}
       \le C\norm U_{\mathrm{HS}}\norm V_{\mathrm{HS}}.
 \label{mat:weak-sandwich}
\end{equation}
For any compatible square matrix $D$,
\begin{equation}
 \left\|\abs{\tr\bigl(D(\one-zB)^{-1}\bigr)}\right\|_{L^{1,\infty}}
       \le C\norm D_1.
 \label{mat:weak-trace}
\end{equation}
Here and below $\norm D_1$ denotes the trace norm. In particular, every
moment of order $0<p<1$ is bounded by $C_p$ times the corresponding
right-hand scale to the power $p$.
\end{lemma}

\begin{proof}
We first prove a weak estimate for the conjugate Poisson kernel and then
represent the matrix resolvent by a positive matrix-valued measure.
Let $\nu$ be a finite measure with values in a Hilbert space,
and let $m_0$ be its total variation. Uniformly in $0\le\rho<1$,
\begin{equation}
 \norm{\norm{Q_\rho*\nu}}_{L^{1,\infty}}\le C m_0,
 \qquad
 \int_\T\norm{P_\rho*\nu}\,\dd m_\T\le m_0.
 \label{mat:hilbert-kernel}
\end{equation}
The Poisson assertion follows from positivity and Fubini. If $m_0=0$,
both assertions are immediate. Otherwise, for the conjugate assertion,
identify the circle with a normalized interval of length one and fix a
level $t\ge2m_0$. Choose maximal dyadic intervals
$I$ for which $\abs\nu(I)>t\abs I$, where $\abs\nu$ denotes total
variation. They are disjoint and have total length at most $m_0/t$.
Their dyadic parents are not selected, so
$\abs\nu(I)\le2t\abs I$.

Replace $\nu$ on each selected interval by its vector average. The
resulting good measure has a density $g$ with
\[
 \norm g_{L^\infty}\le2t,\qquad
 \norm g_{L^1}\le m_0,\qquad
 \norm g_{L^2}^2\le2tm_0.
\]
These assertions include a possibly singular original measure. The
variation measure restricted to the complement of the selected
intervals is bounded by $t\abs J$ on every dyadic interval $J$: if $J$
lies inside a selected interval the restriction vanishes there, and
otherwise $J$ cannot have variation greater than $t\abs J$. Dyadic
coverings extend this domination to Borel sets, proving absolute
continuity and the density bound. The Fourier multipliers of $Q_\rho$
have modulus at most one,
and Hilbert-space Parseval therefore gives
$\norm{Q_\rho*g}_{L^2}^2\le2tm_0$.

The bad pieces
\[
 \nu_I=\nu|_I-\frac{\nu(I)}{\abs I}\,\one_I\,\dd m_\T
\]
have zero total mass and total variations summing to at most $2m_0$.
Interpret a concentric triple interval of length at least one as the
whole circle. Outside the triple interval $3I$, cancellation and the
uniform kernel derivative bound
\[
 \abs{Q_\rho'(u)}\le\frac C{\dist(u,0)^2}
\]
give
\[
 \int_{\T\setminus3I}\norm{Q_\rho*\nu_I}\,\dd m_\T
       \le C\abs{\nu_I}(I).
\]
Indeed, subtract the kernel evaluated at the center of $I$, bound its
difference by $C\abs I/\dist(u,I)^2$, and integrate away from $3I$.
For completeness, in angular coordinates and $\rho>0$,
\[
 \abs{Q_\rho'(u)}
 \le\frac{2\rho}{\abs{1-\rho e^{iu}}^2}
 \le\frac1{2\sin^2(u/2)}
 \le\frac C{\dist(u,2\pi\Z)^2};
\]
the case $\rho=0$ is trivial. Rescaling the circle only changes the
absolute constant, so the bound is uniform as $\rho$ tends to one. The union
of the triple intervals has length at most $3m_0/t$. Chebyshev applied
to the good $L^2$ bound and the bad $L^1$ bound outside this union gives
$m_\T\{\norm{Q_\rho*\nu}>Ct\}\le Cm_0/t$. Levels below $2m_0$ are
trivial. This proves \Cref{mat:hilbert-kernel}.

Now define the analytic matrix function
\[
 F_1(z)=2(\one-zB)^{-1}-\one.
\]
It satisfies $F_1(0)=\one$ and
\[
 \operatorname{Re}F_1(z)
 = (\one-\bar zB^*)^{-1}
       (\one-\abs z^2B^*B)(\one-zB)^{-1}\ge0.
\]
For $r<r'<1$, the positive matrix-valued measure
\[
 \dd M_{r'}(t)=\operatorname{Re}F_1(r'e^{it})\,\dd m_\T(t)
\]
has total mass $\one$. Fourier series show that on the circle of radius
$r$, $F_1$ is its Poisson integral plus $i$ times its conjugate Poisson
integral, with kernel radius $r/r'$. There is no additional constant
imaginary part because $F_1(0)=\one$.

The measure $U M_{r'} V$, considered in the Hilbert space of matrices
with Hilbert--Schmidt norm, has variation at most
$\norm U_{\mathrm{HS}}\norm V_{\mathrm{HS}}$. To verify this for a
finite measurable partition, write $M_j$ for its positive matrix
masses. Then
\begin{align*}
 \sum_j\norm{UM_jV}_{\mathrm{HS}}
 &\le\sum_j\norm{UM_j^{1/2}}_{\mathrm{HS}}
                \norm{M_j^{1/2}V}_{\mathrm{HS}}\\
 &\le\left(\sum_j\norm{UM_j^{1/2}}_{\mathrm{HS}}^2\right)^{1/2}
       \left(\sum_j\norm{M_j^{1/2}V}_{\mathrm{HS}}^2\right)^{1/2}
  =\norm U_{\mathrm{HS}}\norm V_{\mathrm{HS}}.
\end{align*}
Taking the supremum over partitions proves the claim. Applying
\Cref{mat:hilbert-kernel} proves the required weak estimate for $UF_1V$.
The added constant $UV$ in
$U(\one-zB)^{-1}V=(UF_1(z)V+UV)/2$ obeys the same scale bound.

For the trace assertion, factor $D=XY$ by singular values so that
$\norm X_{\mathrm{HS}}\norm Y_{\mathrm{HS}}=\norm D_1$. For any
positive partition masses $M_j$ of total mass $\one$,
\[
 \sum_j\abs{\tr(DM_j)}
 \le\sum_j\norm{YM_j^{1/2}}_{\mathrm{HS}}
              \norm{M_j^{1/2}X}_{\mathrm{HS}}
 \le\norm D_1.
\]
Thus the scalar measure $\tr(DM_{r'})$ has variation at most the trace
norm of $D$, and the same argument applies.

All estimates are uniform on inner circles. Rationality gives radial
limits outside finitely many boundary points. Applying Fatou to level
sets, first at a slightly smaller level and then increasing it, gives
\Cref{mat:weak-sandwich,mat:weak-trace} on the boundary. Finally, a tail
bound $\PP\{X>t\}\le\min(1,A/t)$ gives
\[
 \E X^p=p\int_0^\infty t^{p-1}\PP\{X>t\}\,\dd t
          \le C_p A^p,
 \qquad 0<p<1.
\]
\end{proof}

We record a useful extension of the same proof. If a rational analytic
$d\times d$ matrix function $F$ has nonnegative Hermitian real part in
the disk, its representing positive measures have total mass
$\operatorname{Re}F(0)$. Their Hilbert--Schmidt variation is at most
$\tr(\operatorname{Re}F(0))\le d\norm{F(0)}$. Restoring the constant
$i\operatorname{Im}F(0)$, the preceding argument gives
\begin{equation}
 \left\|\norm{F(e^{it})}\right\|_{L^{1,\infty}}
       \le C d\norm{F(0)}.
 \label{mat:positive-real-extension}
\end{equation}
We need only the resulting polynomial dependence on dimension and the
norm at the center.

\subsection{Determinant phases and artificial absorptions}

To average inverse traces over a real diagonal shift, we integrate the
phase of the determinant. The estimate must also allow artificial
absorptions that vary with that shift.

\begin{lemma}[Phase integration]\label{mat:phase}
Let $D(s)$ be a continuously differentiable $d\times d$ matrix on a
bounded real interval, with $-\operatorname{Im}D(s)>0$. Its determinant
has a continuous phase $\Phi$ such that
\begin{equation}
 -d\pi<\Phi(s)<0,\qquad
 \Phi'(s)=\operatorname{Im}\tr\bigl(D'(s)D(s)^{-1}\bigr).
 \label{mat:phase-derivative}
\end{equation}
Suppose, in particular, that
\[
 D'(s)=P_J-i\tau\diag_J(a_i'(s)),\qquad
 0\le\tau\le1,\qquad \abs{a_i'(s)}^2\le C a_i(s),
\]
where $P_J$ is the coordinate projection onto $J$, and that
$-\operatorname{Im}D(s)\ge\tau\diag_J(a_i(s))$. With $G=D^{-1}$,
\begin{equation}
 \Phi'(s)\ge\tfrac12\tr_J\operatorname{Im}G(s)-C'd.
 \label{mat:phase-lower}
\end{equation}
Consequently, on an interval $[a,b]$,
\begin{equation}
 \int_a^b\tr_J\operatorname{Im}D(s)^{-1}\,\dd s
                  \le C''d(1+b-a).
 \label{mat:phase-integral}
\end{equation}
The integral bound remains valid for passive matrices at their
almost-everywhere existing inverses, by strict regularization.
\end{lemma}

\begin{proof}
Every eigenvalue of a strictly passive matrix lies in the open lower
half-plane: test the imaginary part of the quadratic form on an
eigenvector. Use the branch of the logarithm whose arguments there lie
in $(-\pi,0)$. The sum of the eigenvalue arguments, counted with
algebraic multiplicity, defines a continuous phase of the determinant
in $(-d\pi,0)$. The logarithmic derivative of the determinant gives
\Cref{mat:phase-derivative}; no differentiability of individual
eigenvalues is needed.

The absorption identity gives, for every $i\in J$,
\[
 (\operatorname{Im}G)_{ii}\ge\tau a_i\abs{G_{ii}}^2.
\]
Thus the derivative contribution of an artificial diagonal satisfies
\[
 \tau\abs{a_i'G_{ii}}
       \le C^{1/2}\sqrt\tau\sqrt{(\operatorname{Im}G)_{ii}}.
\]
This is also valid when $a_i=0$, since then $a_i'=0$. Taking imaginary
parts of the trace derivative and applying Young's inequality yields
\[
 \Phi'
 \ge\sum_{i\in J}(\operatorname{Im}G)_{ii}
       -C^{1/2}\sqrt\tau\sum_{i\in J}
                       \sqrt{(\operatorname{Im}G)_{ii}}
 \ge\tfrac12\tr_J\operatorname{Im}G-C'd.
\]
Integration, followed by $\Phi(b)-\Phi(a)\le d\pi$, proves
\Cref{mat:phase-integral}. For merely passive $D$, apply the same
argument to $D-i\zeta\one$, $\zeta>0$. The bounds are independent of
$\zeta$; where $D^{-1}$ exists, the inverses converge as $\zeta$ tends
to zero, and Fatou proves the stated extension.
\end{proof}

Two variants will be used without changing this argument. If the real
derivative is a Hermitian matrix $K(s)\ge cP_J$ for a fixed $c>0$, its
contribution is $\tr(K\operatorname{Im}G)\ge c\tr_J\operatorname{Im}G$.
Thus \cref{mat:phase-integral} holds with constants also depending on
$c$. Additional real diagonal derivatives can instead be retained explicitly
in \Cref{mat:phase-derivative} and estimated using available absorption
bounds. This is how the calibrated root derivative will be handled.

\subsection{The artificial probe profile}
\label{mat:probes}

Separate from the cell ports, we allow at most one artificial probe at
each site. A prescribed \emph{mask} chooses among absence, strength
$\tau$, and strength $\tau\vartheta(v_i)$, where $v_i$ is always the
true site potential, including when the real reference value differs
from it. Masks do not depend on the potentials. They may be fixed after
conditioning on spatial offsets or other independent prescriptions.
Zero strength means absence.

Here $\vartheta:[-h,h]\to[0,1]$ is fixed, of class $C^{1,1}$, polynomial
on finitely many pieces, positive on $(-h,h)$, and identically one on
a plateau containing $[-h/2,h/2]$ and the compact interior potential
settings of the cell construction. Such a profile can be chosen after
those compact settings are fixed and before the later scale and width
parameters are selected. To construct it, choose $a<h$ large enough
that the required compact set lies in $[-a,a]$, and use
$3t^2-2t^3$ on each endpoint transition, with
$t=(v+h)/(h-a)$ on $[-h,-a]$ and $t=(h-v)/(h-a)$ on $[a,h]$; put
$\vartheta=1$ on $[-a,a]$. The endpoint values are zero, the derivatives
match at the joins, and
\begin{equation}
 \abs{\vartheta'(v)}^2\le C\vartheta(v),\qquad
 \vartheta(v)\asymp\dist(v,\{-h,h\})^2
 \quad\hbox{near the endpoints}.
 \label{mat:profile-properties}
\end{equation}
Constants may depend on the fixed profile. In particular,
$\vartheta^{-q}$ is integrable on $[-h,h]$ for every $q<1/2$.

For estimates at a single cell of size $L$, artificial strengths will
obey $0\le\tau\le e^{-c_*L}$ with a sufficiently large fixed $c_*>0$.
The screening estimates will also allow $0\le\tau\le1$. All spatial
masks remain arbitrary subject to the independence convention above.
